\documentclass[11pt]{amsart}

\usepackage[T1]{fontenc}
\usepackage{lmodern}
\usepackage{microtype}
\usepackage{amsmath,amssymb,amsthm}
\usepackage{booktabs}
\usepackage[hidelinks]{hyperref}

\newtheorem{theorem}{Theorem}
\newcommand{\RM}{\operatorname{RM}}

\title[An $\RM(5,2,2)$ graph]
{A Mixed Radial Moore Graph with Parameters $(5,2,2)$}

\date{}

\subjclass[2020]{05C12, 05C20, 05C35}
\keywords{mixed graph, radial Moore graph, Moore bound,
degree--diameter problem, finite certificate}

\hypersetup{
  pdftitle={A Mixed Radial Moore Graph with Parameters (5,2,2)}
}

\begin{document}

\begin{abstract}
We give an explicit totally $(5,2)$-regular mixed graph on
$52=M(5,2,2)$ vertices with radius $2$ and diameter $3$. Hence an
$\RM(5,2,2)$ mixed radial Moore graph exists, settling a case marked
open in Table~3 of Ceresuela, L\'opez, and Chemisana and still marked
open in Table~4.2 of Abiad, Ceresuela, and L\'opez; among the many
parameter pairs left unresolved in the latter table, $(r,z)=(5,2)$ is
the one of least Moore bound, $M(5,2,2)=52$. The graph has unique central
vertex $0$, total status $5741$, and status $1$-norm $N_1=801$.
No optimality assertion for $N_1$ is made.
\end{abstract}

\maketitle

\section{The construction}

A mixed graph has undirected edges and directed arcs. Distances are
source-to-target distances: edges may be traversed in either direction,
whereas arcs may be traversed only forward. The center $C(G)$ is the set
of vertices of minimum out-eccentricity. A mixed graph is totally
$(r,z)$-regular if every vertex has undirected degree $r$, indegree $z$,
and outdegree $z$. Following \cite{CLC2023}, an $\RM(r,z,k)$ graph is a
totally $(r,z)$-regular mixed graph of radius $k$, diameter $k+1$, and
order equal to the mixed Moore bound $M(r,z,k)$.

For $k=2$, the Moore-tree count is
\begin{equation}\label{eq:moore-bound}
M(r,z,2)
 =1+(r+z)+r(r-1+z)+z(r+z)
 =(r+z)^2+z+1.
\end{equation}
Thus $M(5,2,2)=52$. The existence question for this parameter pair is
marked unresolved in Table~3 of \cite{CLC2023} and Table~4.2 of
\cite{ACL2026}. For $k=2$,
Ceresuela, L\'opez, and Chemisana constructed
$\RM(r,1,2)$ with $r\ge1$ and $\RM(1,z,2)$ with $z\ge1$
\cite{CLC2023}, and Abiad, Ceresuela, and L\'opez settled the six
further cases $(r,z)=(2,2),(2,3),(3,2),(3,3),(4,2),(2,4)$ by integer
programming \cite{ACL2026}. The result below gives a
positive answer to the $(5,2,2)$ case of Question~1 in \cite{CLC2023}.

For a vertex $v$, let
\[
s(v)=\sum_{u\in V(G)}d(v,u)
\]
be its status. The hypothetical mixed Moore distance layers for these
parameters have sizes $1,7,44$, so the corresponding vertex status is
\begin{equation}\label{eq:moore-status}
\sigma_{5,2,2}=7+2\cdot44=95.
\end{equation}
We therefore use
\[
N_1(G)=\sum_{v\in V(G)}\lvert s(v)-95\rvert,
\]
in agreement with the status norm of \cite{CLC2023}.

\begin{theorem}\label{thm:main}
Let $G$ be the mixed graph on $V(G)=\{0,1,\ldots,51\}$ whose undirected
neighborhoods $U(v)$ and outgoing-arc heads $A^+(v)$ are listed in the
adjacency certificate below. Then $G$ is an $\RM(5,2,2)$ graph. Moreover,
\[
C(G)=\{0\},\qquad
\sum_{v\in V(G)}s(v)=5741,\qquad
N_1(G)=801.
\]
\end{theorem}

\begin{proof}
The certificate has no loops, directed digons, or edge--arc overlaps.
Undirected adjacency is symmetric, every row contains five undirected
neighbors and two outgoing arcs, and every vertex appears as an arc head
exactly twice. Thus $G$ is totally $(5,2)$-regular, with $130$ edges and
$104$ arcs.

The vertices at distance one from $0$ are $\{1,\ldots,7\}$. After the
parent edge is omitted for vertices $1,\ldots,5$, their admissible
second-layer descendants, together with those of $6$ and $7$, are
respectively
\[
\begin{gathered}
\{8,\ldots,13\},\ \{14,\ldots,19\},\ \{20,\ldots,25\},
\ \{26,\ldots,31\},\ \{32,\ldots,37\},\\
\{38,\ldots,44\},\ \{45,\ldots,51\}.
\end{gathered}
\]
These sets are disjoint and exhaust $V(G)\setminus\{0,\ldots,7\}$.
Hence the distance layers from $0$ have sizes $1,7,44$, and
$\operatorname{ecc}^+(0)=2$.

Breadth-first search from all $52$ sources in the out-neighborhood
relation $U(v)\cup A^+(v)$ gives
\[
\#\{(u,v):d(u,v)=i\}=
\begin{cases}
52,&i=0,\\
364,&i=1,\\
1487,&i=2,\\
801,&i=3,
\end{cases}
\]
and
\[
\operatorname{ecc}^+(0)=2,
\qquad
\operatorname{ecc}^+(v)=3\quad(1\le v\le51).
\]
Consequently $G$ has radius $2$, diameter $3$, and center $\{0\}$.
Together with \eqref{eq:moore-bound}, this proves that
$G\in\RM(5,2,2)$.

Every vertex has seven targets at distance one. If $q(v)$ is the number
of targets at distance three from $v$, then the remaining $44-q(v)$
targets are at distance two, and hence
\[
s(v)=7+2(44-q(v))+3q(v)=95+q(v).
\]
Since $\sum_v q(v)=801$, we obtain
\[
\sum_v s(v)=52\cdot95+801=5741
\]
and, by \eqref{eq:moore-status},
\[
N_1(G)=\sum_v\lvert s(v)-95\rvert
      =\sum_v q(v)
      =801.
\]
\end{proof}

\paragraph{Verification.}
The adjacency certificate below is self-contained: the structural
conditions, the all-pairs distances, the eccentricities, the statuses,
and $N_1$ are all finite computations on the $52$ printed rows, and we
performed them by computer directly from that table. No ancillary file
is needed to reproduce the checks.

\medskip
\noindent\textbf{Adjacency certificate.}
For each vertex $v$, the following table gives the undirected neighbors
$U(v)$ and the heads $A^+(v)$ of its outgoing arcs.

\begin{center}
\fontsize{6.8}{7.15}\selectfont
\setlength{\tabcolsep}{3.2pt}
\renewcommand{\arraystretch}{0.96}
\begin{tabular}{rll@{\qquad}rll}
\toprule
$v$ & $U(v)$ & $A^+(v)$ & $v$ & $U(v)$ & $A^+(v)$ \\
\midrule
0 & $1,2,3,4,5$ & $6,7$ & 26 & $4,15,22,30,41$ & $31,38$ \\
1 & $0,8,9,10,11$ & $12,13$ & 27 & $4,15,24,36,40$ & $13,22$ \\
2 & $0,14,15,16,17$ & $18,19$ & 28 & $4,19,20,39,51$ & $5,34$ \\
3 & $0,20,21,22,23$ & $24,25$ & 29 & $4,22,32,33,51$ & $2,44$ \\
4 & $0,26,27,28,29$ & $30,31$ & 30 & $14,21,26,39,51$ & $20,42$ \\
5 & $0,32,33,34,35$ & $36,37$ & 31 & $10,32,37,43,44$ & $8,34$ \\
6 & $38,39,40,41,42$ & $43,44$ & 32 & $5,11,29,31,47$ & $9,46$ \\
7 & $45,46,47,48,49$ & $50,51$ & 33 & $5,9,18,22,29$ & $15,19$ \\
8 & $1,36,43,48,50$ & $14,47$ & 34 & $5,13,16,19,39$ & $17,22$ \\
9 & $1,25,33,35,43$ & $16,50$ & 35 & $5,9,14,46,50$ & $18,27$ \\
10 & $1,21,23,31,44$ & $4,33$ & 36 & $8,23,25,27,49$ & $10,12$ \\
11 & $1,12,15,32,40$ & $23,35$ & 37 & $31,38,42,45,46$ & $16,30$ \\
12 & $11,14,16,44,46$ & $32,51$ & 38 & $6,24,37,43,49$ & $0,29$ \\
13 & $16,25,34,46,50$ & $41,45$ & 39 & $6,28,30,34,51$ & $27,48$ \\
14 & $2,12,30,35,45$ & $26,32$ & 40 & $6,11,24,27,50$ & $9,29$ \\
15 & $2,11,23,26,27$ & $17,40$ & 41 & $6,17,22,26,48$ & $21,23$ \\
16 & $2,12,13,34,44$ & $15,41$ & 42 & $6,18,37,47,49$ & $8,25$ \\
17 & $2,19,20,41,47$ & $24,45$ & 43 & $8,9,25,31,38$ & $4,48$ \\
18 & $19,21,33,42,48$ & $1,5$ & 44 & $10,12,16,31,45$ & $43,49$ \\
19 & $17,18,28,34,50$ & $14,42$ & 45 & $7,14,37,44,48$ & $10,11$ \\
20 & $3,17,21,28,47$ & $33,39$ & 46 & $7,12,13,35,37$ & $3,6$ \\
21 & $3,10,18,20,30$ & $1,7$ & 47 & $7,17,20,32,42$ & $3,36$ \\
22 & $3,26,29,33,41$ & $11,28$ & 48 & $7,8,18,41,45$ & $28,35$ \\
23 & $3,10,15,24,36$ & $0,26$ & 49 & $7,25,36,38,42$ & $20,47$ \\
24 & $23,27,38,40,51$ & $2,37$ & 50 & $8,13,19,35,40$ & $21,46$ \\
25 & $9,13,36,43,49$ & $38,39$ & 51 & $24,28,29,30,39$ & $40,49$ \\
\bottomrule
\end{tabular}
\end{center}

\begin{thebibliography}{9}
\small

\bibitem{CLC2023}
J.~M.~Ceresuela, N.~L\'opez, and D.~Chemisana,
\emph{On mixed radial Moore graphs of diameter 3},
Discrete Math. \textbf{346} (2023), Article 113525.
\href{https://doi.org/10.1016/j.disc.2023.113525}
{doi:10.1016/j.disc.2023.113525}.

\bibitem{ACL2026}
A.~Abiad, J.~M.~Ceresuela, and N.~L\'opez,
\emph{A note on integer programming methods for mixed radial Moore graphs},
Discrete Math. Lett. \textbf{17} (2026), 87--92.
\href{https://doi.org/10.47443/dml.2026.107}
{doi:10.47443/dml.2026.107}.

\end{thebibliography}

\end{document}